%% Advanced Robotics (Taylor & Francis) submission -- wrist-orientation-fix honest diagnostic paper
%% Template class verified 2026-07-19 (WebSearch): `interact` is the real, standard Taylor & Francis
%% LaTeX class used across their journals, including Advanced Robotics specifically.
%% NOT verified: this journal's specific fee/OA policy and review-turnaround claims -- check the
%% actual "Instructions for Authors" page on tandfonline.com directly before submitting (automated
%% fetch attempts here returned 403 -- their site blocks non-browser access, so this could not be
%% confirmed from this environment). Also NOT compilable in this environment (verified 2026-08-02:
%% `interact.cls` is not installed here) -- compile in an environment with the class available
%% before submitting; this pass was checked manually for balanced braces/escaping only.
%%
%% STATUS (2026-08-02): REFRAMED, per explicit instruction, around this project's own critical
%% ablation finding (`tango_robot/piper_robosuite/README.md`, "DECISIVE FINDING" 2026-07-21 and its
%% n=152 rescale 2026-07-22): CR-CFM (the learned flow-matching policy) adds no measurable value
%% beyond an independent, training-free wrist-orientation-selection fix, on every object tested, and
%% underperforms it outright on a second object. The paper's headline is now the fix itself -- real,
%% statistically significant at adequate power, mechanistically explained, and shown to generalize
%% (or not) in a predictable way -- with the CR-CFM investigation reported as the honest diagnostic
%% process that found it, not as a validated learned-control contribution. This matches the
%% "Embracing Negative Results" framing (arXiv:2406.03980) the underlying investigation already used
%% internally. All numbers in this file are sourced directly from
%% `tango_robot/piper_robosuite/README.md`'s dated 2026-07-21/22 entries and
%% `cr_cfm/IMPROVEMENT_PLAN.md` (not independently re-verified from raw trial logs by this pass -- if
%% a number here and one in those source files ever disagree, the source files win).
%%
%% Still open before this is submission-ready:
%%   1. Methodology's architecture subsections (Translation-Invariant Conditioning, RHC, Sub-Segment
%%      Augmentation) remain brief comment-only descriptions of the system under diagnosis, not full
%%      prose -- lower priority now that CR-CFM is not the contribution, but should still be written.
%%   2. Repeat-based (not single-run) evaluation is mandatory throughout -- single-run-per-trial was
%%      found to have a real ~12.5% instability rate (Discussion, "Evaluation Protocol Reliability").
%%   3. The "1005-vs-1006 micro-trajectory" investigation is CANCELLED (verified 2026-07-19: a
%%      run-to-run nondeterminism artifact, not a real phenomenon) and does not appear anywhere in
%%      this reframed version -- if a future edit ever reintroduces angle-sparsity-counter-example
%%      language, check it against this note first.
%% See `tango_robot/piper_robosuite/README.md`'s CR-CFM entries for the full, dated history this
%% paper is drawn from -- read the LATEST entries first, since earlier entries in that file describe
%% the now-superseded CR-CFM-as-contribution framing.

\documentclass[]{interact}

\usepackage{epstopdf}
\usepackage[caption=false]{subfig}
\usepackage[numbers,sort&compress]{natbib}
\usepackage{booktabs}
\usepackage{graphicx}
\usepackage{hyperref}
\usepackage{cleveref}

\begin{document}

\title{When a Simple Kinematic Fix Outperforms a Learned Policy: An Honest Diagnostic Evaluation of Flow Matching for Contact-Robust Robotic Descent}

\author{
\name{Author A}
\affil{Department, University}
}

\maketitle

\begin{abstract}
We investigate whether a resource-constrained (49K-parameter, single-GPU) flow-matching policy
(CR-CFM) can replace hand-designed execution-control rules for contact-robust descent in robotic
grasping. An extensive six-stage diagnostic process -- addressing run-to-run reproducibility, a
rejected over-regularization hypothesis, receding-horizon control (RHC) instability, and two
distinct training-distribution conditioning gaps -- appeared to show CR-CFM significantly beating a
plain interpolation baseline (81.2\% vs.\ 59.4\%, McNemar exact $p=0.0156$, $n=32$). A critical
ablation the original comparison had not run overturns that conclusion: the comparison had silently
confounded CR-CFM with an independent, training-free wrist-orientation-selection fix that prevents a
specific joint (wrist-roll) from pinning against its hardware limit during descent. Comparing
baseline+fix against CR-CFM+fix directly yields a perfect tie ($26/32$, zero discordant pairs),
confirmed at five times the sample size ($n=152$: $75.7\%$ vs.\ $72.4\%$, 17 discordant pairs split
11-to-6 toward the baseline, $p=0.33$) and explained architecturally by a perturbation test showing
CR-CFM's closed loop reacts only to its own execution drift, never to the environment. CR-CFM
contributes no measurable value beyond the fix on any tested object, and underperforms it outright
on a second object (Pear: $75\%$ vs.\ $62.5\%$). The fix itself is this paper's real, validated
contribution: statistically significant at adequate power ($n=152$, $p=0.027$, $+7.2$ percentage
points), mechanistically explained (a joint6-pinning correlation, $p=1.8\times10^{-5}$), and
object-dependent in a way we show is predictable in advance via a cheap kinematic-only proxy,
without executing a single physical trial. We report the CR-CFM investigation itself as a rigorous,
honest diagnostic process that found a real, deployable fix -- not as a validated learned-control
contribution.
\end{abstract}

\begin{keywords}
Imitation Learning; Flow Matching; Ablation Analysis; Failure Analysis; Kinematic Constraints; Grasp Orientation Selection
\end{keywords}

% =============================================================================
\section{Introduction}
\label{sec:intro}

Edge-deployed manipulation needs small models; small models attempting generative trajectory
synthesis (flow matching / diffusion) face distribution-shift and dynamic-instability failure modes
not visible at large-model scale. We built CR-CFM, a 49K-parameter conditional flow-matching policy
with receding-horizon control, specifically to solve one such failure mode: brittle open-loop
descent trajectories during contact-robust grasping on a Piper 6-DoF arm in simulation. Six
sequential, empirically-verified diagnostic fixes -- addressing run-to-run reproducibility, a
rejected over-regularization hypothesis, RHC instability, and two distinct training-distribution
conditioning gaps -- eventually produced a policy that significantly outperformed a plain,
non-learned interpolation baseline on our primary test object (Cracker: $81.2\%$ vs.\ $59.4\%$,
McNemar exact $p=0.0156$, $n=32$), the cleanest positive result of the underlying project.

This paper is not that result. Before treating it as a validated contribution, we ran a critical
ablation the original comparison had not: isolating whether the win was actually attributable to
CR-CFM's learned, closed-loop correction, or to an independent grasp-orientation-selection mechanism
(\texttt{wrist\_friendly\_orientation}) that had been silently bundled into the CR-CFM arm but not
the baseline arm. The mechanism is simple and training-free -- it prevents the wrist-roll joint from
being driven toward its hardware limit during descent, a failure mode we trace mechanistically to a
specific kinematic correlation ($p=1.8\times10^{-5}$). Comparing baseline+fix against CR-CFM+fix
directly, controlling for this confound, reveals that CR-CFM adds no measurable value beyond the
fix: a perfect tie at $n=32$, confirmed -- not merely failed to be refuted -- at $n=152$, with an
architectural explanation for why: CR-CFM's closed loop reacts only to its own execution drift,
never to the environment, so nothing in our evaluation protocol ever gave it a chance to earn its
added complexity.

We report this reversal with the same rigor as the original positive claim. The wrist-orientation
fix is real, statistically significant at adequate power, mechanistically grounded, and -- unlike
CR-CFM -- requires no training data, no model, and no real-time inference. It is also
object-dependent, and we show that dependence is predictable in advance, per object, from cheap
kinematic geometry alone. CR-CFM is not this paper's contribution; the diagnostic process that found
the fix, and the fix itself, are.

% =============================================================================
\section{Methodology and Closed-Loop Architecture}
% Framed as the system under diagnosis, not the paper's proposed contribution -- \S4 shows why.

\subsection{Baseline Pipeline (Interpolation / \texttt{descend\_refresh})}
% Plain, non-learned interpolation between the current pose and the IK-solved descend target;
% zero training investment, the reference point every other condition is compared against.

\subsection{The Wrist-Orientation-Selection Fix (\texttt{wrist\_friendly\_orientation})}
This paper's actual validated contribution. At grasp-orientation selection time, choose the
candidate orientation that keeps the wrist-roll joint (\texttt{robot0\_joint6}) away from its
hardware limit, rather than the naive nearest-orientation choice. The mechanism is simple,
training-free, and model-independent: it requires no learned component and applies identically to
any downstream descend-execution method. It is motivated by a diagnosed failure mode -- naive
orientation selection drives joint6 toward pinning against its limit during descent on
geometrically demanding objects (Cracker's naive pinning rate: $34.4\%$), stalling the PD
controller's ability to track the planned trajectory. It is statistically confirmed at the
mechanism level (joint6-pinning correlation with trial failure, $p=1.8\times10^{-5}$) and at the
outcome level, pooled across ranges ($n=152$: fix $73.0\%$ vs.\ baseline $65.8\%$, McNemar exact
$p=0.027$, $+7.2$ percentage points; \cref{tab:wristfix_effect}).

\subsection{Translation-Invariant 6-DoF Conditioning Space}
% CR-CFM's object-relative pose conditioning, retained as system description.

\subsection{Receding-Horizon Control (RHC) Defense Against Open-Loop Instability}
% CR-CFM's RHC loop, retained as system description; \S5.1 shows it re-reads only the arm's own
% state each iteration, never the object's position or a re-solved target.

\subsection{Sub-Segment Training Augmentation}
% CR-CFM's training-distribution augmentation for descend sub-trajectories, retained as system
% description.

% =============================================================================
\section{Experimental Design and Ablation Setup}
% N=30 scale-up protocol on the original Cracker Gate-1 comparison; fresh, untouched probe range
% for final validation (both earlier ranges are tuning-contaminated across the six sequential
% diagnostic decisions -- state this explicitly, do not reuse them for headline numbers).

\subsection{The Critical Ablation: Isolating the Wrist-Fix from CR-CFM}
The original Gate-1 comparison (\cref{sec:intro}) compared plain baseline (no wrist-fix) against
CR-CFM+wrist-fix -- but the wrist-fix is orthogonal to, and independent of, whichever
descend-execution method is used underneath it. This ablation instead compares baseline+wrist-fix
against CR-CFM+wrist-fix directly, using the same 32-trial paired protocol as Gate~1, then rescaled
to $n=152$ (matching the wrist-fix's own confirmatory power level, following a pre-registered
decision rule set before rescaling: if the $n=32$ tie is an underpowered near-miss, rescaling should
reveal a real CR-CFM-specific effect; if the tie is real, rescaling should not manufacture one).

\subsection{Cross-Object Generalization Protocol}
Two follow-up protocols test whether the wrist-fix's benefit generalizes beyond Cracker. First, a
direct isolation on Pear (plain baseline vs.\ baseline+wrist-fix, no CR-CFM in either arm, 8 trials,
3-repeat majority vote per trial) mirrors the Cracker ablation's own isolation logic. Second, a
cheap, execution-free kinematic predictor -- does an object's naive grasp orientation drive joint6
toward its hardware limit, measured via IK solves alone (20 trials/object, no physics execution
beyond transit-high and approach) -- tests whether per-object benefit is predictable in advance,
rather than requiring an expensive execution-based trial for every candidate object.

\subsection{Comparative Evaluation}

\begin{table}[ht]
\tbl{Critical ablation isolating CR-CFM's contribution from the wrist-orientation fix (Cracker,
paired, 3-repeat majority vote per trial).}
{\begin{tabular}{lccc} \toprule
Scale & Baseline+Fix & CR-CFM+Fix & McNemar (exact) \\ \midrule
$n=32$  & $26/32$ ($81.2\%$)  & $26/32$ ($81.2\%$)  & tie, 0 discordant pairs \\
$n=152$ & $115/152$ ($75.7\%$) & $110/152$ ($72.4\%$) & $p=0.33$; 17 discordant (11 baseline, 6 CR-CFM) \\
\bottomrule
\end{tabular}}
\label{tab:critical_ablation}
\end{table}

\begin{table}[ht]
\tbl{The wrist-orientation fix's own effect, isolated from CR-CFM.}
{\begin{tabular}{lccc} \toprule
Object ($n$) & Baseline & Baseline+Fix & McNemar (exact) \\ \midrule
Pooled ($n=152$) & $100/152$ ($65.8\%$) & $111/152$ ($73.0\%$) & $p=0.027$ \\
Pear ($n=8$)     & $6/8$ ($75\%$)       & $6/8$ ($75\%$)       & tie, 0 discordant pairs \\
\bottomrule
\end{tabular}}
\label{tab:wristfix_effect}
\end{table}

\begin{table}[ht]
\tbl{Cheap kinematic (IK-only) predictor of wrist-fix benefit: predicted joint6-pinning rate by
object (20 trials/object, no physics execution beyond transit-high and approach).}
{\begin{tabular}{lc} \toprule
Object & Predicted pinning rate \\ \midrule
Cracker & $6/20$ ($30\%$) \\
Pear & $1/20$ ($5\%$) \\
Mustard & $1/20$ ($5\%$) \\
\bottomrule
\end{tabular}}
\label{tab:pinning_predictor}
\end{table}

% =============================================================================
\section{Discussion}

\subsection{Why CR-CFM Adds Nothing: An Architectural Explanation, Not Just an Empirical One}
Building on the critical ablation's null result (\cref{tab:critical_ablation}), a perturbation test
injecting a one-time joint-space disturbance ($\pm0.1$ rad, uniform) partway
through descent produces identical degradation under both conditions (baseline $3/8$, $37.5\%$;
CR-CFM $3/8$, $37.5\%$ -- still fully concordant, zero discordant pairs). The root cause: the
descend target is IK-solved once before descent starts, for both conditions; CR-CFM's RHC loop
re-reads only the arm's own joint state each iteration, never the object's position or a re-solved
target. Its closed loop closes around its own execution, not around the environment. This is why no
evaluation protocol in this study -- including the perturbation test -- could have given CR-CFM a
chance to earn its added complexity: nothing here tests the one thing a re-solving, object-aware
closed loop would need to demonstrate.

\subsection{Object-Dependence of the Wrist-Fix, and Why It Is Predictable in Advance}
On Pear, the wrist-fix alone shows zero measurable effect ($6/8$ vs.\ $6/8$), mechanistically
confirmed via direct joint6 instrumentation across all 16 trial-condition combinations: zero show
pinning ($|{\rm joint6}| $ within $0.05$ rad of the limit), with the largest magnitude observed still
more than $0.5$ rad from the limit. This is not a failure of the fix -- Pear's grasp geometry (a
small, round object with approach angles already concentrated in a narrow arc) simply never drives
the wrist into the region Cracker's flat, wide box geometry does. CR-CFM+fix underperforms plain
baseline+fix outright on Pear ($75\%$ vs.\ $62.5\%$), reinforcing that CR-CFM's added complexity
buys nothing even where the fix itself has less to contribute. This object-dependence motivated the
cheap, execution-free predictor in \cref{tab:pinning_predictor}, validated to reproduce Cracker's
independently-known pinning rate exactly ($11/32$, trial-for-trial identical to the original
diagnostic) and to correctly predict Mustard's benefit as null before running a single physical
trial -- turning an expensive, per-object empirical question into a cheap geometric one.

\subsection{Evaluation Protocol Reliability: Single-Run Instability}
Approximately $12.5\%$ ($2/16$) of single-run trial classifications disagreed across otherwise-
identical repeats, concentrated in one probe range and absent in the other -- a genuine limitation
of single-run evaluation for this class of system, not an artifact of any individual trial. This
motivated the repeat-based, majority-vote protocol used throughout this paper as a methodological
necessity, worth stating plainly as standard practice for future evaluation in this area rather than
trusting any single execution as ground truth.

% =============================================================================
\section{Conclusion and Future Work}

We set out to validate a learned flow-matching policy for contact-robust descent and, through the
same rigorous diagnostic discipline that motivated building it, found instead that the entire
measured improvement is attributable to a much simpler, training-free mechanism: a
wrist-orientation-selection fix that prevents one joint from pinning against its hardware limit.
This is not a result to report quietly -- it is the paper's actual contribution. The fix is
statistically significant at adequate power ($n=152$, $p=0.027$), mechanistically grounded
($p=1.8\times10^{-5}$), requires no training data, no model, and no real-time inference, and is
shown to be predictable in advance, per object, via a cheap kinematic proxy rather than expensive
execution. CR-CFM itself -- the 49K-parameter policy, its receding-horizon control loop, and every
CR-CFM-specific mechanism developed across six diagnostic stages -- contributes no measurable value
beyond the fix on any object tested, and underperforms it outright on a second object. We report the
CR-CFM investigation as a rigorous, transparent diagnostic process that led to a real, useful,
immediately deployable finding, not as a validated learned-control algorithm. Future work should
treat the wrist-orientation-selection fix, not CR-CFM, as the system worth deploying or extending --
for example, generalizing the underlying joint-limit-avoidance principle to other kinematically
constrained joints, or validating the cheap IK-only predictor against a broader object set than the
seven currently available in this simulation registry.

\end{document}
